% Portada de C1: un mismo modelo al que se llega por tres vías (aplicación, Claude Code y API).
\documentclass[tikz,border=0pt]{standalone}
\usepackage{resonancia}
\begin{document}
\begin{tikzpicture}
  \clip[rounded corners=6pt] (0,0) rectangle (12,3.6);
  \shade[top color=rGris!22, bottom color=white] (0,0) rectangle (12,3.6);
  % modelo
  \draw[trazo, color=rVioleta, fill=white, rounded corners=8pt] (4.3,0.5) rectangle (7.7,1.5);
  \node[etiqueta, color=rVioleta] at (6,1.0) {modelo};
  % tres vías
  \foreach \x/\t/\c in {1.6/aplicación/rAzul, 6/Claude Code/rVerde, 10.4/API/rNaranja} {
    \draw[trazo, color=\c, fill=white, rounded corners=6pt] (\x-1.4,2.2) rectangle (\x+1.4,3.2);
    \node[etiqueta, color=\c] at (\x,2.7) {\t};
  }
  \draw[trazo, color=rAzul, -{Stealth}] (1.6,2.2) -- (4.3,1.3);
  \draw[trazo, color=rVerde, -{Stealth}] (6,2.2) -- (6,1.5);
  \draw[trazo, color=rNaranja, -{Stealth}] (10.4,2.2) -- (7.7,1.3);
  \node[etiqueta, anchor=north, color=rGris] at (1.6,2.05) {chat, proyectos, archivos};
  \node[etiqueta, anchor=north, color=rGris] at (10.4,2.05) {tu propio código};
  \node[etiqueta, anchor=north, color=rGris] at (6,0.35) {agente en tu terminal};
\end{tikzpicture}
\end{document}
